% !TeX root = ../main.tex
% Adaptasi struktur contoh LaTeX ITERA v2026.1; lihat catatan sumber.
% Kolom tanda tangan, foto, dan tanggal sengaja belum diisi.
\begingroup
\setstretch{1}
\setlength{\parskip}{6pt}
\setlength{\tabcolsep}{0pt}
\halamanAwal{LEMBAR PENGESAHAN}

Tugas Akhir berjudul ``\judulIndonesia'' disusun oleh \namaPenulis, NIM \nomorInduk, pada Program Studi \namaProgramStudi, Fakultas \namaFakultas, \namaInstitusi.

\noindent
\begin{minipage}[t]{0.69\linewidth}
  \vspace{0pt}
  \namaLokasi, \tanggalPengesahan\par
  Penulis,\par
  \vspace{1.5cm}
  \namaPenulis\par
  NIM \nomorInduk
\end{minipage}\hfill
\begin{minipage}[t]{0.25\linewidth}
  \vspace{0pt}
  \centering
  \setlength{\fboxsep}{0pt}
  \fbox{\parbox[c][4cm][c]{3cm}{\centering Pasfoto\\3 $\times$ 4 cm}}
\end{minipage}

\vspace{0.4cm}
\begin{center}
  Diperiksa dan disetujui oleh,
\end{center}

\noindent
\begin{tabular}{@{}p{0.7\linewidth}>{\centering\arraybackslash}p{0.3\linewidth}@{}}
  \textbf{Pembimbing} & \textbf{Tanda tangan}\\[6pt]
  1. \pembimbingSatu & \\
  \hspace*{1em}\identitasPembimbingSatu & \dotfill\\[12pt]
  2. \pembimbingDua & \\
  \hspace*{1em}\identitasPembimbingDua & \dotfill\\[12pt]
  \textbf{Penguji} & \\[6pt]
  1. \pengujiSatu & \\
  \hspace*{1em}\identitasPengujiSatu & \dotfill\\[12pt]
  2. \pengujiDua & \\
  \hspace*{1em}\identitasPengujiDua & \dotfill
\end{tabular}

\vspace{0.5cm}
\begin{center}
  Disahkan oleh,\\
  Koordinator Program Studi \namaProgramStudi\\
  Fakultas \namaFakultas\\
  \namaInstitusi\par
  \vspace{1.4cm}
  \koordinatorProdi\\
  \identitasKoordinator
\end{center}
\clearpage
\endgroup
